% META: {"id": "1SPE-SECDEG-CO-024", "chapitre": "1SPE-SECOND-DEGRE", "type_objet": "corrige", "exercice_ref": "1SPE-SECDEG-EX-024", "capacites_codes": ["C3", "C4"], "sources_inspiration": [], "status": "generated"}
% BEGIN-VERIFY
% from sympy import symbols, solve, simplify, Rational
% x = symbols('x')
% # Q1 : (x^2 - 4)(x - 3) = 0
% f = (x**2 - 4)*(x - 3)
% sols = sorted(solve(f, x))
% assert sols == [-2, 2, 3]
% # Q2 : signe de g(x) = x^2 - 4 = (x-2)(x+2)
% g = x**2 - 4
% assert g.subs(x, 0) < 0
% assert g.subs(x, 3) > 0
% assert g.subs(x, -3) > 0
% # Q3 : h(x) = 2x^2 + 3x - 5
% h = 2*x**2 + 3*x - 5
% delta_h = 3**2 - 4*2*(-5)
% assert delta_h == 49
% sols_h = sorted(solve(h, x))
% assert sols_h == [Rational(-5, 2), 1]
% assert simplify(h - (2*x + 5)*(x - 1)) == 0
% # Q4 : h >= 0 sur ]-inf ; -5/2] U [1 ; +inf[
% assert h.subs(x, 0) < 0
% assert h.subs(x, 2) > 0
% assert h.subs(x, -3) > 0
% END-VERIFY
\begin{corrige}{1SPE-SECDEG-EX-024}

\textbf{1. Résolution de $(x^2 - 4)(x - 3) = 0$.}

\commentaireMarge{Règle du produit nul : $AB = 0 \iff A = 0$ ou $B = 0$.}

On factorise $x^2 - 4 = (x-2)(x+2)$, donc :
\[(x-2)(x+2)(x-3) = 0.\]

Par la règle du produit nul :
\[x - 2 = 0 \text{ ou } x + 2 = 0 \text{ ou } x - 3 = 0,\]
ce qui donne $x = 2$, $x = -2$ ou $x = 3$.

L'ensemble des solutions est $\mathcal{S} = \{-2\,;\,2\,;\,3\}$.

\medskip
\textbf{2. Tableau de signes de $g(x) = x^2 - 4$.}

\commentaireMarge{$g(x) = (x-2)(x+2)$, racines $-2$ et $2$, $a = 1 > 0$.}

$g(x) = x^2 - 4 = (x-2)(x+2)$. Les racines sont $-2$ et $2$.

\medskip
\begin{center}
\begin{tabular}{|c|ccccccc|}
\hline
$x$ & $-\infty$ & & $-2$ & & $2$ & & $+\infty$ \\
\hline
$g(x)$ & & $+$ & $0$ & $-$ & $0$ & $+$ & \\
\hline
\end{tabular}
\end{center}

\medskip
\textbf{3. Étude de $h(x) = 2x^2 + 3x - 5$.}

\commentaireMarge{$a=2$, $b=3$, $c=-5$.}

\textit{a) Discriminant et racines.}

\[\Delta = 3^2 - 4 \times 2 \times (-5) = 9 + 40 = 49.\]

\[x_1 = \frac{-3 - 7}{4} = \frac{-10}{4} = -\frac{5}{2} \qquad \text{et} \qquad x_2 = \frac{-3 + 7}{4} = \frac{4}{4} = 1.\]

\textit{b) Factorisation.}

\[h(x) = 2\!\left(x + \frac{5}{2}\right)\!(x - 1) = (2x + 5)(x - 1).\]

\textit{Vérification :} $(2x+5)(x-1) = 2x^2 - 2x + 5x - 5 = 2x^2 + 3x - 5$. \hfill $\square$

\textit{c) Tableau de signes de $h$.}

\commentaireMarge{$a = 2 > 0$ : signe $+$ à l'extérieur des racines.}

\medskip
\begin{center}
\begin{tabular}{|c|ccccccc|}
\hline
$x$ & $-\infty$ & & $-\dfrac{5}{2}$ & & $1$ & & $+\infty$ \\
\hline
$h(x)$ & & $+$ & $0$ & $-$ & $0$ & $+$ & \\
\hline
\end{tabular}
\end{center}

\medskip
\textbf{4. Solutions de $2x^2 + 3x - 5 \geqslant 0$.}

\commentaireMarge{Lire les intervalles où $h \geqslant 0$ dans le tableau.}

D'après le tableau de signes, $h(x) \geqslant 0$ sur $\left]-\infty\,;\,-\dfrac{5}{2}\right] \cup \left[1\,;\,+\infty\right[$.

\baremeIndicatif{Q1 : 2 pts (calculer, raisonner) ; Q2 : 2 pts (raisonner) ; Q3a-b : 2 pts (calculer) ; Q3c : 1 pt (raisonner) ; Q4 : 1 pt (raisonner)}
\end{corrige}
